% GENERATED FILE. Edit canonical lessons, then run npm run generate:books.
% canonical-source-hash: fnv1a64:be49baca46dd16cf
% canonical-lessons: GU-C255-varta, GU-C255-kavita, GU-C255-navalkatha, GU-C255-natak, GU-C255-philm

\chapter{Stories, poems and films}
\label{ch:gu-stories-poems-and-films}
\hlchaptermodality{255}

\begin{chapteropening}
\textbf{By the end of this chapter:} \emph{I can say which stories I heard and which films and plays I saw.}

Complete the last lesson of chapter 255: I can say which stories I heard and which films and plays I saw.
\end{chapteropening}

\section[\texorpdfstring{\gu{વાર્તા} (vārtā)}{વાર્તા (vārtā)}]{\gu{વાર્તા} (vārtā) — story}

\label{lesson:GU-C255-varta}



\begin{quote}
Before the new one: say the Gujarati for flute, then the Gujarati for tabla drum.
\end{quote}

\subsection*{You'll want to know: \gu{વાર્તા}}
\textbf{\gu{વાર્તા}} — \emph{vārtā} — \textquotedblleft{}story\textquotedblright{}.

\begin{quote}
\textbf{\gu{દાદીએ} \gu{અમને} \gu{વાર્તા} \gu{કહી}.} — \emph{dādīe amne vārtā kahī} — \textquotedblleft{}Grandmother told us a story.\textquotedblright{}
\end{quote}

\begin{cousinweb}
\textbf{\gu{વાર્તા}} (\emph{vārtā}) comes from Sanskrit \emph{vārttā}, \textquotedblleft{}livelihood; tidings\textquotedblright{}, in a learned shape, a semi-tatsama (CDIAL 11564). CDIAL lists the inherited \emph{vāt}, \textquotedblleft{}news\textquotedblright{}, from the same word.
\end{cousinweb}

\begin{grammarlens}[title={What you've built}]
One more word for everyday life at A2.
\end{grammarlens}

\subsection*{Your turn}
\begin{itemize}
\raggedright
  \item \emph{Say it:} \emph{vārtā}
  \item \emph{Say it:} \emph{vārtā}, once more
  \item \emph{Say it:} \emph{dādīe amne vārtā kahī} — Grandmother told us a story.
  \item \emph{Recall:} say the Gujarati for flute, then the Gujarati for tabla drum, then say \emph{vārtā} again
\end{itemize}

\begin{tcolorbox}[breakable,colback=teal!4,colframe=teal!35!black,title={Before you move on}]
What does \gu{વાર્તા} mean? (\textquotedblleft{}Story\textquotedblright{}.) Say it once more.
\end{tcolorbox}
\section[\texorpdfstring{\gu{કવિતા} (kavitā)}{કવિતા (kavitā)}]{\gu{કવિતા} (kavitā) — poem, poetry}

\label{lesson:GU-C255-kavita}



\begin{quote}
Before the new one: say the Gujarati for tabla drum, then the Gujarati for story.
\end{quote}

\subsection*{You'll want to know: \gu{કવિતા}}
\textbf{\gu{કવિતા}} — \emph{kavitā} — \textquotedblleft{}poem, poetry\textquotedblright{}.

\begin{quote}
\textbf{\gu{મેં} \gu{શાળામાં} \gu{એક} \gu{કવિતા} \gu{વાંચી}.} — \emph{mẽ śāḷāmā̃ ek kavitā vā̃cī} — \textquotedblleft{}I read a poem at school.\textquotedblright{}
\end{quote}

\begin{cousinweb}
\textbf{\gu{કવિતા}} (\emph{kavitā}) is Sanskrit \emph{kavitā}, \textquotedblleft{}poetry\textquotedblright{}, taken over whole as a tatsama. It is built on \emph{kaví}, \textquotedblleft{}wise man, poet\textquotedblright{} (CDIAL 2964).
\end{cousinweb}

\begin{grammarlens}[title={What you've built}]
One more word for everyday life at A2.
\end{grammarlens}

\subsection*{Your turn}
\begin{itemize}
\raggedright
  \item \emph{Say it:} \emph{kavitā}
  \item \emph{Say it:} \emph{kavitā}, once more
  \item \emph{Say it:} \emph{mẽ śāḷāmā̃ ek kavitā vā̃cī} — I read a poem at school.
  \item \emph{Recall:} say the Gujarati for tabla drum, then the Gujarati for story, then say \emph{kavitā} again
\end{itemize}

\begin{tcolorbox}[breakable,colback=teal!4,colframe=teal!35!black,title={Before you move on}]
What does \gu{કવિતા} mean? (\textquotedblleft{}Poem, poetry\textquotedblright{}.) Say it once more.
\end{tcolorbox}
\section[\texorpdfstring{\gu{નવલકથા} (navalkathā)}{નવલકથા (navalkathā)}]{\gu{નવલકથા} (navalkathā) — novel}

\label{lesson:GU-C255-navalkatha}



\begin{quote}
Before the new one: say the Gujarati for story, then the Gujarati for poem.
\end{quote}

\subsection*{You'll want to know: \gu{નવલકથા}}
\textbf{\gu{નવલકથા}} — \emph{navalkathā} — \textquotedblleft{}novel\textquotedblright{}.

\begin{quote}
\textbf{\gu{તેણે} \gu{મોટી} \gu{નવલકથા} \gu{વાંચી}.} — \emph{teṇe moṭī navalkathā vā̃cī} — \textquotedblleft{}He read a long novel.\textquotedblright{}
\end{quote}

\begin{cousinweb}
\textbf{\gu{નવલકથા}} (\emph{navalkathā}) is a modern learned compound built from \emph{naval}, \textquotedblleft{}new, wonderful\textquotedblright{}, and Sanskrit \emph{kathā}, \textquotedblleft{}conversation; story\textquotedblright{} (CDIAL 2704).
\end{cousinweb}

\begin{grammarlens}[title={What you've built}]
One more word for everyday life at A2.
\end{grammarlens}

\subsection*{Your turn}
\begin{itemize}
\raggedright
  \item \emph{Say it:} \emph{navalkathā}
  \item \emph{Say it:} \emph{navalkathā}, once more
  \item \emph{Say it:} \emph{teṇe moṭī navalkathā vā̃cī} — He read a long novel.
  \item \emph{Recall:} say the Gujarati for story, then the Gujarati for poem, then say \emph{navalkathā} again
\end{itemize}

\begin{tcolorbox}[breakable,colback=teal!4,colframe=teal!35!black,title={Before you move on}]
What does \gu{નવલકથા} mean? (\textquotedblleft{}Novel\textquotedblright{}.) Say it once more.
\end{tcolorbox}
\section[\texorpdfstring{\gu{નાટક} (nāṭak)}{નાટક (nāṭak)}]{\gu{નાટક} (nāṭak) — drama, play}

\label{lesson:GU-C255-natak}



\begin{quote}
Before the new one: say the Gujarati for poem, then the Gujarati for novel.
\end{quote}

\subsection*{You'll want to know: \gu{નાટક}}
\textbf{\gu{નાટક}} — \emph{nāṭak} — \textquotedblleft{}drama, play\textquotedblright{}.

\begin{quote}
\textbf{\gu{અમે} \gu{શનિવારે} \gu{નાટક} \gu{જોયું}.} — \emph{ame śanivāre nāṭak joyũ} — \textquotedblleft{}We saw a play on Saturday.\textquotedblright{}
\end{quote}

\begin{cousinweb}
\textbf{\gu{નાટક}} (\emph{nāṭak}) is Sanskrit \emph{nāṭaka}, \textquotedblleft{}dancer, actor\textquotedblright{} and later \textquotedblleft{}drama\textquotedblright{}, taken over as a tatsama (CDIAL 7045).
\end{cousinweb}

\begin{grammarlens}[title={What you've built}]
One more word for everyday life at A2.
\end{grammarlens}

\subsection*{Your turn}
\begin{itemize}
\raggedright
  \item \emph{Say it:} \emph{nāṭak}
  \item \emph{Say it:} \emph{nāṭak}, once more
  \item \emph{Say it:} \emph{ame śanivāre nāṭak joyũ} — We saw a play on Saturday.
  \item \emph{Recall:} say the Gujarati for poem, then the Gujarati for novel, then say \emph{nāṭak} again
\end{itemize}

\begin{tcolorbox}[breakable,colback=teal!4,colframe=teal!35!black,title={Before you move on}]
What does \gu{નાટક} mean? (\textquotedblleft{}Drama, play\textquotedblright{}.) Say it once more.
\end{tcolorbox}
\section[\texorpdfstring{\gu{ફિલ્મ} (philm)}{ફિલ્મ (philm)}]{\gu{ફિલ્મ} (philm) — film, movie}

\label{lesson:GU-C255-philm}



\begin{quote}
Before the new one: say the Gujarati for novel, then the Gujarati for drama.
\end{quote}

\subsection*{You'll want to know: \gu{ફિલ્મ}}
\textbf{\gu{ફિલ્મ}} — \emph{philm} — \textquotedblleft{}film, movie\textquotedblright{}.

\begin{quote}
\textbf{\gu{ગઈકાલે} \gu{મેં} \gu{નવી} \gu{ફિલ્મ} \gu{જોઈ}.} — \emph{gaīkāle mẽ navī philm joī} — \textquotedblleft{}Yesterday I saw a new film.\textquotedblright{}
\end{quote}

\begin{cousinweb}
\textbf{\gu{ફિલ્મ}} (\emph{philm}) is a loan from English \emph{film}.
\end{cousinweb}

\begin{grammarlens}[title={What you've built}]
One more word for everyday life at A2.
\end{grammarlens}

\subsection*{Your turn}
\begin{itemize}
\raggedright
  \item \emph{Say it:} \emph{philm}
  \item \emph{Say it:} \emph{philm}, once more
  \item \emph{Say it:} \emph{gaīkāle mẽ navī philm joī} — Yesterday I saw a new film.
  \item \emph{Recall:} say the Gujarati for novel, then the Gujarati for drama, then say \emph{philm} again
\end{itemize}

\begin{tcolorbox}[breakable,colback=teal!4,colframe=teal!35!black,title={Before you move on}]
What does \gu{ફિલ્મ} mean? (\textquotedblleft{}Film, movie\textquotedblright{}.) Say it once more.
\end{tcolorbox}
